% !TEX root =  Lecture21.tex


%%%%%%%%%%%%%%%%%%%%%%%%%%%%%%%%%%%%
% Recap/Agenda
%%%%%%%%%%%%%%%%%%%%%%%%%%%%%%%%%%%%
\begin{frame}
\frametitle{Module 4: Bayesian Statistics}
    \begin{itemize}
        \item \hl{Previously:} how to \emph{read} a posterior: point estimates, credible intervals, hypothesis tests, model comparison, prediction, and predictive checks.
        \item \hl{This time:} the \hl{regression model} enters the Bayesian frame, the last and biggest model of the course.
          \begin{itemize}
              \item the model: a slope, an intercept, and noise, each with \hl{its own prior};
              \item \hl{choosing those priors}, and \hl{checking them before the data} by simulating lines from the prior alone;
              \item then Lecture 20 applies unchanged: a \hl{credible interval} for the slope, a \hl{prediction} for a new penguin.
          \end{itemize}
        \item \hl{Reading:} Bayes Rules!\ Chapter 9.
        \item \hl{Homework:}
    \end{itemize}
\end{frame}


%%%%%%%%%%%%%%%%%%%%%%%%%%%%%%%%%%%%
% Learning objectives:
%%%%%%%%%%%%%%%%%%%%%%%%%%%%%%%%%%%%
\begin{frame}
    \frametitle{Lecture Objectives}
    \goals{%
    \begin{itemize}
    \item write down a \hl{Bayesian linear regression}: the same data model as OLS, plus a prior on each unknown;
    \item choose a prior for each parameter and say whether it is \hl{vague}, \hl{weakly informative}, or \hl{informative}: and why that choice is defensible;
    \item run a \hl{prior predictive check}: simulate regression lines from the prior alone and judge whether they are plausible;
    \item read a \hl{credible interval for the slope} off the posterior, and say how the prior moved it relative to OLS;
    \item predict a new case, distinguishing the \hl{line} $E[y \mid x]$ from a \hl{new observation}: the Bayesian answer to Lecture 15's CI-versus-PI.
    \end{itemize}}
\end{frame}


%%%%%%%%%%%%%%%%%%%%%%%%%%%%%%%%%%%%
\begin{frame}
    \frametitle{Lecture Objectives}
    \objectives{%
    \begin{itemize}
    \item \textbf{(M4, LO6)} Bayesian Linear Regression
    \end{itemize}}
    \vspace{0.3em}
    {\small Building on, and revisiting:}
    \objectives{%
    \begin{itemize}
    \item \textbf{(M4, LO4)} Inference and Prediction: Lecture 20's credible intervals and posterior prediction, now applied to a regression
    \item \textbf{(M4, LO5)} Markov chain Monte Carlo: Lecture 19's sampler, used here to fit three parameters at once
    \item \textbf{(M3)} Regression inference: the OLS fit, bootstrap, and prediction intervals we are about to reinterpret
    \end{itemize}}
\end{frame}
